\chapter{Two weak complementary events force the C2 radial certificate}
\label{component:p534f57e1da6f}
\noindent\textbf{Component }\path{p534f57e1da6f}\par
{\small\noindent Source: \path{CORE/contact/chart_free_c2.tex}\par}
{\small\noindent SHA256: \path{534f57e1da6f10b1f7e4c37aff830d3ad6bea003ec6350fabd1bdf5cedd3e7b7}\par}
{\small\noindent Inventory use: Only exact physical witnesses, contact amplitude estimate, relative coalescence, scaled stationary signs and CT8 radial source entry; its all-tests calls are resolved by the actual radial/ordinary bodies listed here\par}
\medskip
\section{Two weak complementary events force the C2 radial certificate}
\label{p534f57e1da6f:sec:chart-free-c2}

Here weak means infinitesimal after division by the original positive
signed coefficient \(d\). All source calculations take place in one real
closed coefficient field. Testing fields remain independent.

\begin{theorem}[C2 entry without a contact-chart hypothesis]
\label{p534f57e1da6f:thm:chart-free-c2}
Let the physical canonical parent have \(\lambda=\beta=0\), \(d>0\),
and nonnegative original event coefficients. Suppose it has an actual
finite noncommuting stationary zero contact, and
\[
                    C/d,\ \alpha/d\quad\text{are infinitesimal}.
\]
Then its full parent is tensor SOS at every independent real or complex
projective test. No hypothesis on the other four event weights, on the
contact-angle ratio, or on the signs of the source gaps is imposed.

More precisely, normalize \(d=1\). After local real sign gauges and at
most the exact C2 reflection,
the actual contact enters the positive rational C2 source with
\[
 c\text{ bounded},\qquad a,b,t\longrightarrow1,\qquad
 r,h,g,j>0\text{ infinitesimal}.
\]
In particular all four matching weights are unbounded and both
complementary weights are strictly positive on this finite-contact domain.
\end{theorem}

The proof has three ingredients: a finite physical witness makes the
matching weights large; the quadratic CHSH excess controls the contact
amplitudes at the scale of the actual angles; stationarity fixes their
signs. This last step establishes the source chart rather than assuming it.

\subsection{The physical bound at the correct angular scale}
For \(u,v>0\), put \(D_u=1+u^2,D_v=1+v^2\), and let \(M(u,v)\) be the
real symmetric matrix of the signed core
\[
H=2\{R_D^2+D_u(vY-Z)^2+D_v(X+uZ)^2-D_uD_vZ^2\},
\]
where
\[
R_m=vW-X+uvY-uZ,\qquad R_D=uW+uvX-Y-vZ.
\]
Write \(L=D_uD_v\). Direct matrix multiplication gives
\begin{equation}
(M-LI)^2=L^2I+4uvLJ,\qquad
J=\begin{pmatrix}0&0&0&1\\0&0&-1&0\\0&-1&0&0\\1&0&0&0\end{pmatrix},
\qquad J^2=I.                                      \label{p534f57e1da6f:eq:cf-core-square}
\end{equation}
For a unit real state \(\psi\), Cauchy--Schwarz therefore gives
\[
\psi^TM\psi\ge L-\sqrt{L^2+4uvL}\ge-2uv.
\]
The last inequality follows by squaring the positive quantities
\(L+2uv\) and \(\sqrt{L^2+4uvL}\); their squared difference is
\(4u^2v^2\). Homogeneity proves
\begin{equation}
                         H\ge-2uv\,\|\psi\|^2.        \label{p534f57e1da6f:eq:cf-core-bound}
\end{equation}
This is an exact ordered-field inequality, not a truncated expansion.
It is the usual quadratic CHSH excess in the present normalization.

At a normalized zero contact, every event term is nonnegative. Define
the four matching amplitudes, in the fixed order \((n,A,m,D)\), by
\[
\mathcal L=(\sqrt{D_uD_v}X,\sqrt{D_uD_v}Y,R_m,R_D).
\]
Equation~\eqref{p534f57e1da6f:eq:cf-core-bound} gives, for each matching event,
\begin{equation}
                   e_i\mathcal L_i^2\le uv/2.         \label{p534f57e1da6f:eq:cf-amplitude-bound}
\end{equation}
All square roots and inequalities in this paragraph are coefficient-side.

\subsection{Matching weights and the compact contact geometry}
First every matching weight is unbounded. To see the mechanism for
\(D\), use the coefficient-field experiment
\[
 u=1,\quad v=\sqrt{1+\delta},\quad
 (W,X,Y,Z)=(1,0,0,\sqrt{1+\delta}),\qquad\delta>0.
\]
The \(n,A,m\) packets vanish, and the full eight-event parent is exactly
\[
4D\delta^2-2\delta+8\lambda+4(2+\delta)C
 +8(1+\delta)\alpha+4(2+\delta)(1+\delta)\beta.
\]
If \(D\le B\) for a positive real \(B\), choose the positive real
\(\delta=1/[4(B+1)]\). The first two terms have a strictly negative
real upper bound, whereas the complementary terms are infinitesimal.
This contradicts physicality. The exact even rotations of the event
cycle give the same argument for \(n,A,m\), preserving the sum of the
four complementary weights. The accompanying checker also supplies
the four witnesses directly. Hence \(1/n,1/A,1/m,1/D\) are infinitesimal.

Normalize the actual state to norm one and use positive actual
half-angles, obtainable by local real basis gauges. The weaker
bound \(H/(D_uD_v)\ge-2\) shows that all normalized matching
probabilities vanish in residue. In particular \(X,Y\to0\).
Put
\[
c_u=(1+u^2)^{-1/2},\ s_u=uc_u,\qquad
c_v=(1+v^2)^{-1/2},\ s_v=vc_v.
\]
If \(\sigma=\operatorname{res}(c_us_v)\) and
\(\tau=\operatorname{res}(s_uc_v)\), the two remaining matching
equations read
\[
\sigma\overline W-\tau\overline Z=0,\qquad
\tau\overline W-\sigma\overline Z=0.
\]
At least one of \(\overline W,\overline Z\) is nonzero. Thus
\(\sigma=\tau\ge0\). If positive, the two actual half-angles are
bounded units with the same residue. If zero, both are infinitesimal
or both unbounded, by the two unit-circle identities. The exact C2
reflection sends \((u,v)\) to \((1/v,1/u)\) and exchanges
\((C,\alpha)\) and \((m,D)\), preserving the hypotheses. We may
therefore assume either ordinary bounded-unit angles or two
infinitesimal angles.

\subsection{Relative coalescence when both angles vanish}
Suppose \(u,v\) are infinitesimal. From
\eqref{p534f57e1da6f:eq:cf-amplitude-bound}, all four
\(\mathcal L_i/\sqrt{uv}\) are infinitesimal. In particular
\(X/\sqrt{uv},Y/\sqrt{uv}\to0\). Set \(a=v/u>0\).
Dividing the two other matching amplitudes by \(\sqrt{uv}\) gives
\[
\sqrt a\,W-\frac Z{\sqrt a}\longrightarrow0,\qquad
\frac W{\sqrt a}-\sqrt a\,Z\longrightarrow0.
\]
The discarded terms are \(X/\sqrt{uv},Y/\sqrt{uv}\) and
\(\sqrt{uv}X,\sqrt{uv}Y\), all infinitesimal.
If \(a\) were unbounded, these equations would force both \(W,Z\)
infinitesimal; the same follows if \(a\) is infinitesimal.
Both contradict normalization. Thus \(a\) is a bounded unit.
Taking residues now yields \(a\to1\) and \(W/Z\to1\), with
both \(W,Z\) units. Consequently
\[
                  X/u,\ Y/u,\ R_m/u,\ R_D/u\longrightarrow0.
\]
This is relative coalescence, which the compact-angle argument alone
does not establish.

Normalize \(W=1\) by a pure coefficient factor. Write
\[
x=X/u,\quad y=Y/u,\quad q=v/u,\qquad
(u,x,y,Z,q)\longrightarrow(0,0,0,1,1).
\]
The core divided by \(u^2\) is the polynomial
\begin{align}
\mathcal H:=H/u^2=2\{&
(1-y-qZ+u^2qx)^2+(1+u^2q^2)(x+Z)^2 \notag\\
&+(1+u^2)(u^2q^2y^2-2qyZ-q^2Z^2)\}. \label{p534f57e1da6f:eq:cf-scaled-core}
\end{align}
It has zero residue at this point. The nonnegative alpha packet and
the zero-contact equation give
\[
0\le4\frac{\alpha}{u^2}(1+u^2)(u^2qy-Z)^2
                   \le-\mathcal H.
\]
Its square factor is a positive unit. Hence
\begin{equation}
                         \alpha/u^2\longrightarrow0. \label{p534f57e1da6f:eq:cf-alpha-scale}
\end{equation}
This estimate is derived before the scaled stationary equations are used.

\subsection{Stationarity fixes the source signs}
In the small-angle case introduce the four exact local coordinates
\[
\ell=\left(\sqrt{D_uD_v}\,x,\sqrt{D_uD_v}\,y,
 q-x+u^2qy-Z,\ 1+u^2qx-y-qZ\right).
\]
Their Jacobian with respect to \((x,y,Z,q)\) has residue
\[
J_0=\begin{pmatrix}
1&0&0&0\\0&1&0&0\\-1&0&-1&1\\0&-1&-1&-1
\end{pmatrix},\qquad \det J_0=2.
\]
Thus its inverse and the transformed derivatives are bounded at the
actual contact. Only this pointwise coefficient-field inverse is used.
Direct differentiation of \eqref{p534f57e1da6f:eq:cf-scaled-core} gives
\[
\operatorname{res}(\nabla_\ell\mathcal H)=2(1,-1,-1,1).
\]
The full parent divided by \(u^2\) is exactly
\[
4\sum_i e_i\ell_i^2
 +4C(1+u^2q^2)(1-y)^2
 +4(\alpha/u^2)(1+u^2)(u^2qy-Z)^2+\mathcal H.
\]
The two complementary-gradient terms have zero residue, by weak \(C\)
and \eqref{p534f57e1da6f:eq:cf-alpha-scale}. The original stationary equations imply
stationarity in \(\ell\), with \(u\) fixed. These equations first show
that every \(e_i\ell_i\) is bounded, so its residue is defined. Consequently
\begin{equation}
\operatorname{res}(n\ell_n,A\ell_A,m\ell_m,D\ell_D)
                         =\tfrac14(-1,1,1,-1).        \label{p534f57e1da6f:eq:cf-sign-residue}
\end{equation}
Every matching weight is positive, so \(X<0,Y>0,R_m>0,R_D<0\)
as exact source inequalities.

For ordinary angles the argument is shorter. Normalize \(W=1\);
the actual contact has residue
\((X,Y,Z,v)=(0,0,1,u)\), where \(u\) is a bounded unit.
Use the unscaled \(\mathcal L\) coordinates. Their Jacobian determinant
at that residue is \(2u(1+u^2)^2\), and
\[
\nabla_{\mathcal L}H=2u(1,-1,-1,1)
\]
there. All complementary gradients are bounded times \(C,\alpha\),
hence infinitesimal. Stationarity gives
\(\operatorname{res}(e_i\mathcal L_i)
 =\overline u\,(-1,1,1,-1)/4\), with the same exact signs.

\subsection{The original source and all independent tests}
In either case \(W/Z\to1\), \(v/u\to1\), and \(X/u,Y/u\to0\).
Set, using the original actual contact,
\[
c=u,\quad a=v/u,\quad b=aW/Z,\quad
r=-X/(uZ),\quad t=(uW-Y)/(vZ).
\]
Then \(a,b,t\to1\) and
\[
h=b-a^2t=\frac{aY}{uZ}>0,\qquad
g=\frac{R_m}{uZ}>0,\qquad j=-\frac{R_D}{vZ}>0.
\]
All four \(r,h,g,j\) are infinitesimal, so \(0<r<1\).
The remaining chart boundaries are positive units:
\[
1-c^2h=\frac{Z-vY}{Z}>0,\qquad
b-a^2c^2r=\frac{a(W+vX)}Z>0.
\]
These definitions give exactly
\[
h=b-a^2t,\qquad
g=(1+c^2)b-(1-r+a^2c^2t),\qquad j=1+c^2r-t.
\]
There is no remaining sign-entry assumption.

Finally the stationary coefficient matrix has the nonzero determinant
\[
4^6[-64ac^6rgt\,hj(1-c^2h)(1+c^2)^2(1+a^2c^2)^2\Delta].
\]
The inherited short three-by-three reduction proves this formula.
For \(z=c^2\), its positive expansion is
\[
\Delta=1+(1+z)b+z(1+z)b^2+2z(1+z)br+z^2r^2+2zr
       +a^2z\{(1+z)b+1+zr(2-r)\}\ge1.
\]
The original coefficients
therefore equal the unique positive rational C2 source formulas.
In particular their positive numerator identity forces \(C,\alpha>0\).
If \(c\) is infinitesimal, the complete small-angle radial theorem,
including its outer-test corollary, applies. If \(c\) is a unit,
the inherited ordinary-angle radial corollary applies. The precise
source files are identified in the reproduction paragraph below. These certificates concern
the full parent and arbitrary independent testing fields. Reflect back
when required and restore \(d\) by a pure coefficient square root.
This proves Theorem~\ref{p534f57e1da6f:thm:chart-free-c2}.

The conclusion closes the full two-weak-complementary-event
\emph{finite-contact} C2 branch, not just a preselected rational chart.
It does not assert that every weak-\(C\) source has weak \(\alpha\).
The explicit one-weak escape family in the following supplement shows
why that remaining implication fails.
